\uuid{GFrG}
\titre{ Comptage de paramètres d'un MLP }
\niveau{L3}
\module{Analyse de données}
\chapitre{Réseaux de neurones}
\sousChapitre{Architecture, paramètres}
\theme{Nombre de paramètres, couche cachée, entrée aplatie}
\auteur{Maxime Nguyen}
\datecreate{2026-03-19}
\organisation{AMSCC}
\difficulte{1}

\contenu{
\texte{
  On considère un réseau de neurones entièrement connecté avec une entrée de dimension $3$,
  une couche cachée de $4$ neurones ReLU, et une sortie scalaire.
}
\begin{enumerate}

\item
  \question{Combien de paramètres ce réseau contient-il au total ? Détailler par couche.}
  \reponse{
    \begin{itemize}
      \item Couche 1 (entrée $\to$ cachée) : $3 \times 4 = 12$ poids $+ 4$ biais $= 16$ paramètres.
      \item Couche 2 (cachée $\to$ sortie) : $4 \times 1 = 4$ poids $+ 1$ biais $= 5$ paramètres.
    \end{itemize}
    Total : $16 + 5 = \mathbf{21}$ paramètres.
  }

\item
  \question{Si on double la taille de la couche cachée (8 neurones), combien de paramètres supplémentaires ajoute-t-on ?}
  \reponse{
    Avec $h = 8$ :
    \begin{itemize}
      \item Couche 1 : $3 \times 8 + 8 = 32$ paramètres ($+16$ par rapport à $h=4$).
      \item Couche 2 : $8 \times 1 + 1 = 9$ paramètres ($+4$ par rapport à $h=4$).
    \end{itemize}
    Paramètres supplémentaires : $16 + 4 = \mathbf{20}$.
  }

\item
  \question{Même question pour une image $8 \times 8$ aplatie en entrée avec une couche cachée de $32$ neurones.}
  \reponse{
    L'entrée aplatie est de dimension $64$.
    \begin{itemize}
      \item Couche 1 : $64 \times 32 + 32 = 2080$ paramètres.
      \item Couche 2 : $32 \times 1 + 1 = 33$ paramètres.
    \end{itemize}
    Total : $2080 + 33 = \mathbf{2113}$ paramètres.
  }

\end{enumerate}
}
